%================================================
% CAPÍTULO ARRAYS
%================================================

\chapterimage{chapter_pointers.jpg} % Chapter heading image

\chapter{Arrays}
Quem nunca precisou organizar dados semelhantes num único local? Eu sim, você não? Então quando você, por exemplo, escreve um texto, o que acha que está fazendo? \say{\textit{Organizando ideias?}}. Sim, também! Contudo, o que você está organizando materialmente são caracteres. Vejamos:

\begin {center}
  \begin{bytefield}[bitwidth=1em]{17}
    \bitbox{17}{Autor presunçoso!}\\
    \bitbox{1}{A} \bitbox{1}{u} \bitbox{1}{t} \bitbox{1}{o} \bitbox{1}{r} \bitbox{1}{ }
    \bitbox{1}{p} \bitbox{1}{r} \bitbox{1}{e} \bitbox{1}{s} \bitbox{1}{u} \bitbox{1}{n}
    \bitbox{1}{ç} \bitbox{1}{o} \bitbox{1}{s} \bitbox{1}{o} \bitbox{1}{!}
  \end{bytefield}
\end{center}

A frase \textit{Autor presunçoso!} é a organização de dezessete caracteres um atrás do outro. Isso mesmo, dezessete, pois o espaço também é um dado. Nunca se esqueça que, quando tratamos dados, até um \say{0} ocupa espaço na memória do computador. Engenheiros, acalmem-se, não estou falando de estado negativo num barramento de dados. \textit{All right}?

Segue outra organização de dados. Desta vez de inteiros que serão o resultado do próximo sorteio da mega-sena.

\begin{center}
  \begin{bytefield}[endianness=little,bitwidth=6em]{6}
    \bitbox{1}{9} \bitbox{1}{20} \bitbox{1}{11}
    \bitbox{1}{7} \bitbox{1}{43} \bitbox{1}{31}\\
    % \bitheader{0-5}
  \end{bytefield}
\end{center}

Ambos são exemplos de \textbf{arrays} (também conhecidos como \textit{vetores}) que, em C, são tipos de dados de alocação contíngua derivados dos tipos simples. Em outras palavras, são estruturas de dados homogêneas, compostas de elementos de mesmo tipo, e posicionados em sequência na memória.

\begin{remark}
Em algumas fontes esparsas em idioma português, você, leitor, encontrará referências a arrays como sendo \textit{\textbf{arranjos}}. Ressalte-se que um array, numa tradução mais fiel, é um \textit{\textbf{conjunto}} de dados. Já um arranjo:

  \begin{figure}[!htp]
    \centering
    \includegraphics[scale=0.7, keepaspectratio=true]{arrays_arrange.jpg}
    \caption{Array não é Arranjo}
    \label{fig:arrays_arrange}
  \end{figure}
\end{remark}

\section{Trabalhando com Arrays}
\subsection{Definindo Arrays}

Ao se definir um vetor é necessário determinar o respectivo tamanho.

\begin{ccode}
  int i[6];
  char c[3];
\end{ccode}

O uso do operador de índice [] é que informa ao compilador que as variáveis i e c são do tipo vetor. Já o número entre os colchetes define a quantidade de elementos a serem alocados para o vetor em questão. No código acima, i é definido como um vetor de int com seis elementos e c um vetor de char com três elementos.

\begin{center}
  \begin{bytefield}[endianness=little,bitwidth=6em]{6}
    i[6]\\
    \bitbox{1}{-57} \bitbox{1}{3} \bitbox{1}{791431480}
    \bitbox{1}{0} \bitbox{1}{32767} \bitbox{1}{0}\\
    \bitheader{0-5}\\
  \end{bytefield}
  \begin{bytefield}[endianness=little,bitwidth=2em]{3}
    c[3]\\
    \bitbox{1}{0} \bitbox{1}{-10} \bitbox{1}{49}\\
    \bitheader{0-2}
  \end{bytefield}
\end{center}

Uma boa prática de programação é inicializar os vetores, assim como os tipos simples, pois, tanto em i[6] quanto em c[3], vê-se claramente que os elementos contêm valores aleatórios.

\subsection{Declarando Arrays}

Ao contrário do que se pensa comumente, vetores podem ser declarados sem especificação da quantidade de elementos.

\begin{ccode}
  extern int ext[];
\end{ccode}

\subsection{Inicializando Arrays}

\begin{ccode}
  static int stc[6];
\end{ccode}

\subsection{Acessando Arrays}

\section{Array de Caracteres}

\section{Arrays na Vida Real?!}
\subsection{Imprimir Elementos}
\subsection{Somar Elementos}
\subsection{Inverter o Array}
\subsection{Ordenar o Array}

\section{Arrays como Parâmetros}

\section{Arrays Multidimensionais (Matrizes)}


%------------------------------------------------
